%% kapa: exam=1 round=34 year=2023 period=1 form=A subject=real_estate subject_ko=부동산학원론 seq=14 number=94
%% source: Q-Net 감정평가사 기출문제 (cst00309) artlSeq=5211789 file=raw/1008/5211789/2023년도 제34회 감정평가사 1차시험 1교시.pdf
%% conversion: kordoc 4.24.51 -> tools/gen1.py (auto); review=proofread-sonnet; flags=none
\begin{question}{94}
외부효과에 관한 설명으로 옳은 것은?
\begin{choices}
  \item 외부효과란 거래 당사자가 시장메카니즘을 통하여 상대방에게 미치는 유리하거나 불리한 효과를 말한다.
  \item 부(-)의 외부효과는 의도되지 않은 손해를 주면서 그 대가를 지불하지 않는 외부경제라고 할 수 있다.
  \item 정(+)의 외부효과는 소비에 있어 사회적 편익이 사적 편익보다 큰 결과를 초래한다.
  \item 부(-)의 외부효과에는 보조금 지급이나 조세경감의 정책이 필요하다.
  \item 부(-)의 외부효과는 사회적 최적생산량보다 시장생산량이 적은 과소생산을 초래한다.
\end{choices}
\end{question}
\answer{3}
